\chapter*{Remerciements}
\addcontentsline{toc}{chapter}{Remerciements}
\thispagestyle{plain}

Au terme de ce mini-projet, nous tenons à remercier toutes les personnes qui ont contribué à sa réalisation.

Nos premiers remerciements s'adressent à \textbf{Prof. Maria Zemzami}, qui a proposé et encadré ce projet dans le cadre de la thématique Industrie~4.0. Son cahier des charges --- contexte et objectifs, spécifications fonctionnelles et techniques, aspects organisationnels et humains, livrables, planning et critères d'évaluation --- a structuré notre démarche du début à la fin, et ce rapport en reprend volontairement le plan.

Nous remercions également le corps professoral de l'École Nationale Supérieure d'Arts et Métiers (\ac{ENSAM}) de Rabat, et particulièrement les enseignants de la filière \textit{INDIA\&SD}, pour la qualité de la formation dispensée~: les compétences en programmation, en modélisation et en analyse de données acquises au cours de notre cursus ont été mobilisées à chaque étape du projet.

Enfin, chaque membre de l'équipe remercie les cinq autres pour la rigueur et l'entraide dont chacun a fait preuve~: un jumeau numérique n'existe que si la simulation, la communication, la \ac{3D} et le tableau de bord parlent le même langage, et cet accord s'est construit ensemble.

\vspace{1.5cm}
\begin{flushright}
\textit{A.~\textsc{Issami}, M.~\textsc{Ezbakhe}, H.~\textsc{Benzahra},}\\
\textit{Z.~\textsc{Hsain}, H.~\textsc{Kichchou}, M.~\textsc{Hamdouch}}\\
\textit{2026}
\end{flushright}
